\documentclass[11pt]{article}

\usepackage{nmrdoc}
\usepackage{siunitx}
\usepackage{bm}
\usepackage{tikz}

\sisetup{per-mode=symbol}
\DeclareSIUnit{\angstrom}{\text{\AA}}

\newcommand{\code}[1]{\texttt{#1}}
\setlength{\parindent}{0pt}
\setlength{\parskip}{0.42em}
\emergencystretch=2em

\begin{document}
\NmrTitle{Numerical Addendum: Dispersion Calculator}

The Dispersion calculator is a cutoff sum of closed-form per-vertex London
kernels. The numerical work is the shared ring geometry it consumes, the gates
that decide which ring vertices enter the sum, the CHARMM-form switching
window whose middle branch matches value and first derivative at the joins, and
the shared \(T_0/T_1/T_2\) (trace/3, antisymmetric, symmetric-traceless)
projection used for output. The listings below are
distilled from the source: loops and bookkeeping are trimmed, \code{cfg(...)}
stands for \code{CalculatorConfig::Get(...)}, and the constants, inequalities,
guards, and arithmetic are the implementation's.

\section{Ring vertices and gates}

\code{GeometryResult} fills \code{conf.ring\_geometries} before this
calculator runs. Each ring geometry carries ordered vertex positions
\(\bm v_j\), their mean centre \(\bm c\), an SVD-fitted unit normal, and a mean
centre-to-vertex radius \(R\). The Haigh--Mallion numerical addendum treats the
SVD plane fit in full. Dispersion consumes the vertices for the kernel, the
centre and radius for acceptance, and the normal only for the shared
ring-neighbour coordinate record.

For each atom, the spatial index proposes ring centres within
\(\SI{15.0}{\angstrom}\). That query does not set the dispersion range; the
\(\SI{5.0}{\angstrom}\) range is enforced per vertex. The ring-level
near-field filter uses \code{source\_extent = 2*geom.radius} and
\code{near\_field\_exclusion\_ratio = 0.5}, so the per-atom path accepts only
centre distances greater than \(R\). At equality the ring is omitted. The
topology gate then omits the whole ring if the target atom is a ring vertex or
is covalently bonded to any vertex.

\begin{codebox}
auto nearby_rings =
    spatial.RingsWithinRadius(atom_pos,
        cfg("ring_current_spatial_cutoff"));       // 15.0 A

ctx.distance = (atom_pos - geom.center).norm();
ctx.source_extent = 2.0 * geom.radius;             // ring diameter
if (!filters.AcceptAll(ctx)) continue;             // requires distance > R

if (ring_bonded[ri].count(ai)) continue;           // vertex or bonded neighbour

Mat3 K_ring = Mat3::Zero();
double s_ring = 0.0;
int contacts = 0;

for (size_t vi = 0; vi < ring.atom_indices.size(); ++vi) {
    Vec3 vpos = geom.vertices[vi];
    double r = (atom_pos - vpos).norm();
    auto vertex = ComputeDispVertex(atom_pos, vpos, r);
    if (!vertex.valid) continue;

    K_ring += vertex.K;
    s_ring += vertex.scalar;
    contacts++;
}
\end{codebox}

These gates are omissions. A rejected ring or vertex contributes no tensor, no
scalar contact, and no compensating weight elsewhere in the sum.

\section{One vertex kernel}

For a target atom at \(\bm r_i\) and a ring vertex at \(\bm v_j\), the helper
forms \code{d\_av}, the vector
\(\bm d=\bm r_i-\bm v_j\) from vertex to atom, and uses the caller's
\(r=|\bm d|\). The code does not store \(\hat{\bm d}\), but the unit direction
\(\hat{\bm d}=\bm d/r\) is the angular part of the same arithmetic:
\[
  3\,\frac{d_a d_b}{r^8} - \frac{\delta_{ab}}{r^6}
  =
  \frac{1}{r^6}\left(3\hat d_a\hat d_b-\delta_{ab}\right).
\]
Both terms carry units of \(\si{\angstrom^{-6}}\); the helper applies
\(C_6=1\).

\begin{codebox}
if (r < cfg("singularity_guard_distance")) return result;  // < 0.1 A
if (r > cfg("dispersion_vertex_distance_cutoff")) return result;

double S = DispersionSwitchingFunction(r);
if (S < cfg("dispersion_switching_noise_floor")) return result;

Vec3 d_av = atom_pos - vertex_pos;         // vertex -> atom
double r2 = r * r;
double r6 = r2 * r2 * r2;
double r8 = r6 * r2;

result.scalar = S / r6;

for (int a = 0; a < 3; ++a)
    for (int b = 0; b < 3; ++b) {
        double Iab = (a == b ? 1.0 : 0.0);
        result.K(a,b) = S * (3.0*d_av(a)*d_av(b)/r8 - Iab/r6);
    }

result.valid = true;
\end{codebox}

The tensor part is the symmetric-traceless angular shape
\(3\hat{\bm d}\otimes\hat{\bm d}-\bm I\), scaled by \(S/r^6\). Along
\(\hat{\bm d}\) it has eigenvalue \(2S/r^6\); in either perpendicular
direction it has eigenvalue \(-S/r^6\). Those three eigenvalues sum to zero, so
each accepted vertex contributes no tensor trace and no antisymmetric part.

The scalar contact stored beside the tensor is the isotropic \(S/r^6\) piece.
It is not the trace average of \(\bm K\), which is zero by construction. The
\(-\bm I/r^6\) term inside \(\bm K\) is the subtraction that removes the trace
from \(3\bm d\otimes\bm d/r^8\); the separate \code{scalar} keeps the positive
\(S/r^6\) contact magnitude rather than a trace component.

\section{The switching window}

The switch reads
\(R_s=\SI{4.3}{\angstrom}\) and \(R_c=\SI{5.0}{\angstrom}\) from the shared
TOML file. The code returns full weight through the switch onset, zero weight
at and beyond the cutoff, and the Brooks--CHARMM polynomial between them:
\[
S(r)=
\begin{cases}
1, & r\le R_s,\\[3pt]
\dfrac{(R_c^2-r^2)^2(R_c^2+2r^2-3R_s^2)}
      {(R_c^2-R_s^2)^3}, & R_s<r<R_c,\\[9pt]
0, & r\ge R_c .
\end{cases}
\]

\begin{codebox}
double DispersionSwitchingFunction(double r) {
    double r_switch = cfg("dispersion_switching_onset_distance"); // 4.3 A
    double r_cut    = cfg("dispersion_vertex_distance_cutoff");    // 5.0 A
    if (r <= r_switch) return 1.0;
    if (r >= r_cut)    return 0.0;

    double Rs2 = r_switch * r_switch;
    double Rc2 = r_cut * r_cut;
    double r2  = r * r;
    double num = (Rc2 - r2) * (Rc2 - r2) * (Rc2 + 2.0*r2 - 3.0*Rs2);
    double den = (Rc2 - Rs2) * (Rc2 - Rs2) * (Rc2 - Rs2);
    return num / den;
}
\end{codebox}

The endpoint values come from the factors in the middle branch. At
\(r=R_s\), the numerator equals the denominator, so \(S=1\). At \(r=R_c\), the
double factor \((R_c^2-r^2)^2\) makes \(S=0\). Differentiating the middle
branch gives
\[
  S'(r)=
  \frac{12r(R_c^2-r^2)(R_s^2-r^2)}{(R_c^2-R_s^2)^3},
\]
so \(S'(R_s)=0\) and \(S'(R_c)=0\). Since the outside branches are constant,
the value and first derivative match at both joins. The code does not make
curvature continuous; the second derivative of the middle branch is not
constrained to equal the zero curvature of either constant branch.

\begin{figure}[ht]
\centering
\begin{tikzpicture}[x=1.0cm,y=2.15cm,line width=0.45pt,>=stealth]
  \draw[->] (-0.85,0) -- (5.85,0) node[right] {\scriptsize \(r\)};
  \draw[->] (0,-0.05) -- (0,1.18) node[above] {\scriptsize \(S(r)\)};
  \draw[black!45,densely dotted] (0,1) -- (5.0,1);
  \draw[black!45,densely dotted] (4.3,0) -- (4.3,1.04);
  \draw[black!45,densely dotted] (5.0,0) -- (5.0,1.04);
  \draw (0,1) -- (4.3,1);
  \draw[smooth]
    plot coordinates {
      (4.30,1.000000) (4.35,0.987345) (4.40,0.951194)
      (4.45,0.894447) (4.50,0.820248) (4.55,0.731994)
      (4.60,0.633341) (4.65,0.528212) (4.70,0.420805)
      (4.75,0.315596) (4.80,0.217356) (4.85,0.131149)
      (4.90,0.062345) (4.95,0.016628) (5.00,0.000000)
    };
  \draw (5.0,0) -- (5.72,0);
  \node[below] at (4.3,0) {\scriptsize \(R_s\)};
  \node[below] at (5.0,0) {\scriptsize \(R_c\)};
  \node[left] at (0,1) {\scriptsize \(1\)};
  \node[left] at (0,0) {\scriptsize \(0\)};
\end{tikzpicture}
\caption{The switch is flat at full weight when it leaves the raw kernel at
\(R_s\), and flat again when it reaches zero at \(R_c\). The plotted curve uses
the configured switch for \(R_s=\SI{4.3}{\angstrom}\) and
\(R_c=\SI{5.0}{\angstrom}\).}
\end{figure}

With a hard cutoff, a continuously moving vertex that crosses \(R_c\) would
change the scalar term from approximately \(R_c^{-6}\) to zero in one frame
step, and the tensor term from
\(R_c^{-6}(3\hat{\bm d}\otimes\hat{\bm d}-\bm I)\) to zero. The per-frame sum
would inherit that finite jump from whichever vertex crossed. With this
window, the switched scalar \(S/r^6\) and switched tensor
\(S r^{-6}(3\hat{\bm d}\otimes\hat{\bm d}-\bm I)\) are already zero at
\(R_c\), and their first derivatives with respect to \(r\) are zero there. At
\(R_s\), \(S=1\) and \(S'=0\), so the derivative of the switched kernel matches
the derivative of the raw kernel as the taper begins. The window removes value
and slope discontinuities at the two switch boundaries; it does not remove a
curvature change.

The \(\mathrm C^1\) property belongs to \code{DispersionSwitchingFunction}.
The emitted kernel adds one more gate: \code{ComputeDispVertex} drops any
vertex whose switch value falls below \(10^{-15}\), so near \(R_c\), where
\(S\) has already fallen below that floor, the kernel is set to zero
rather than its taper value. The comparison is strict, so a value exactly at
the floor is kept.

\section{Accumulation and tensor split}

After at least one vertex survives, the ring tensor, scalar contact, and
contact count are stored on the ring-neighbour record. The tensor is also added
to the atom total. Per ring type, the code accumulates the scalar contact sum
and only the five \(T_2\) components of the ring tensor. The per-type arrays
hold \code{kAromaticRingTypeCount = 8} slots, indices \(0\) through \(7\); the
proline pyrrolidine, type index \(8\), is the first saturated ring and falls
outside the \code{ti < kAromaticRingTypeCount} guard, so its dispersion never
reaches a per-type slot.

\begin{codebox}
if (contacts == 0) continue;

ring_neighbour->disp_tensor = K_ring;
ring_neighbour->disp_spherical = SphericalTensor::Decompose(K_ring);
ring_neighbour->disp_scalar = s_ring;
ring_neighbour->disp_contacts = contacts;

int ti = ring.TypeIndexAsInt();
if (ti >= 0 && ti < kAromaticRingTypeCount) {
    ca.per_type_disp_scalar_sum[ti] += s_ring;
    for (int c = 0; c < kT2Components; ++c)
        ca.per_type_disp_T2_sum[ti][c] += ring_neighbour->disp_spherical.T2[c];
}

disp_total += K_ring;
ca.disp_shielding_contribution =
    SphericalTensor::Decompose(disp_total);
\end{codebox}

The exact arithmetic would keep every accepted \(\bm K\) in \(T_2\) only. The
code does not symmetrize or remove the trace after summing, so any floating
residue in \(T_0\) or \(T_1\) remains in \code{disp\_shielding.npy}. The
per-type \(T_2\) array writes only the five stored shape components. The file
\code{disp\_per\_type\_T0.npy} writes per-atom, per-aromatic-type sums of
\(s_{\mathrm{ring}}\), not the tensor trace average.

The shared projection is closed-form arithmetic on the \(3\times3\) tensor
\code{sigma} (here the ring or atom-total dispersion tensor). The trace divided
by three gives \(T_0\), the antisymmetric part gives \(T_1\), and the symmetric
traceless part gives the five \(T_2\) numbers.

\begin{codebox}
st.T0 = sigma.trace() / 3.0;

st.T1[0] = 0.5*(sigma(1,2) - sigma(2,1));
st.T1[1] = 0.5*(sigma(2,0) - sigma(0,2));
st.T1[2] = 0.5*(sigma(0,1) - sigma(1,0));

double Sxx = sigma(0,0) - st.T0;
double Syy = sigma(1,1) - st.T0;
double Szz = sigma(2,2) - st.T0;
double Sxy = 0.5*(sigma(0,1) + sigma(1,0));
double Sxz = 0.5*(sigma(0,2) + sigma(2,0));
double Syz = 0.5*(sigma(1,2) + sigma(2,1));

st.T2[0] = std::sqrt(2.0) * Sxy;
st.T2[1] = std::sqrt(2.0) * Syz;
st.T2[2] = std::sqrt(3.0 / 2.0) * Szz;
st.T2[3] = std::sqrt(2.0) * Sxz;
st.T2[4] = (Sxx - Syy) / std::sqrt(2.0);
\end{codebox}

The \(T_2\) packing is isometric. The \(\sqrt2\) factors on off-diagonal
entries account for the Frobenius double count, since \(S_{xy}\) appears as
both \(S_{xy}\) and \(S_{yx}\). The two diagonal entries
\(\sqrt{3/2}\,S_{zz}\) and \((S_{xx}-S_{yy})/\sqrt2\) carry the remaining two
degrees of freedom after \(S_{xx}+S_{yy}+S_{zz}=0\). A dot product of two
stored \(T_2\) five-vectors is therefore the Frobenius inner product of their
symmetric traceless matrices.

\section{Glossary}

\textbf{London dispersion kernel.} The per-vertex tensor is a signed ellipsoid
that stretches along the vertex-to-atom direction and shrinks equally in the
two transverse directions. One accepted vertex at distance \(r\) contributes
eigenvalues \(2S/r^6,-S/r^6,-S/r^6\), while its separate scalar contact is
\(S/r^6\). Formally,
\(K_{ab}=S(r)r^{-6}(3\hat d_a\hat d_b-\delta_{ab})\), a symmetric traceless
second-order tensor with units \(\si{\angstrom^{-6}}\).

\textbf{\(\mathrm C^1\) switching window.} The switch is a taper that leaves
the kernel unchanged through \(R_s\), lowers it across the cutoff interval, and
reaches zero with a horizontal tangent. In this calculator
\(R_s=\SI{4.3}{\angstrom}\), \(R_c=\SI{5.0}{\angstrom}\), and a vertex at the
exact cutoff is removed because \(S(R_c)=0\) is below the numerical floor.
Formally, the piecewise function has continuous value and first derivative at
\(R_s\) and \(R_c\), with no condition imposed on the second derivative.

\textbf{Dyadic product.} The grid a single direction builds by multiplying each
of its components against every other (so one vector sets both rows and columns,
and the matrix carries that one direction). For a vertex-to-atom unit direction,
\(\hat{\bm d}\otimes\hat{\bm d}\) has entry \(\hat d_a\hat d_b\). Formally, for
column vectors \(\bm u,\bm v\), \((\bm u\otimes\bm v)_{ab}=u_av_b\).

\textbf{Spherical tensor split.} A \(3\times3\) tensor holds three kinds of
content that rotate without mixing --- an orientation-free size (the scalar
trace, \(T_0\)), an axial twist from its antisymmetric part (the three \(T_1\)
components), and a five-number traceless shape for the rest of the
direction-dependence (the five \(T_2\) components). In exact arithmetic, dispersion's tensor sum lives in
the five \(T_2\) components, while its positive \(1/r^6\) contact sum is stored
separately. Formally, it is the closed-form split of a \(3\times3\) tensor's
nine components into dimensions \(1+3+5\), with the \(T_2\) basis normalized to
preserve the Frobenius norm.

\end{document}
